% answers_short.tex —— 2008 年·普通高等学校招生考试·重庆卷·文科 · 简短答案速查
% 由 tools/make_short_answers.py 自 answers.tex 的 \gksolhead/\gkanswer 提取，
% 并与 2008/2008_数学真题_重庆文_答案要点.md 逐题比对后生成，勿手改（改源 MD 后重跑）。

\begin{gknotice}
\textbf{使用说明}：本节只给最终答案便于快速核对，完整推导见解析版。答案为编写者独立推导，\textbf{非官方标准答案}。
\end{gknotice}

\gkpart{选择题}
\begin{center}\begin{tabular}{ccccccccccccc}
\toprule
题号 & 1 & 2 & 3 & 4 & 5 & 6 & 7 & 8 & 9 & 10 & 11 & 12 \\
\midrule
答案 & C & A & C & A & D & D & B & C & B & B & B & C \\
\bottomrule
\end{tabular}\end{center}

\gkpart{填空题}
\begin{itemize}[leftmargin=2.6em]
  \item[\textbf{13.}] $\{2,3\}$。
  \item[\textbf{14.}] $-23$。
  \item[\textbf{15.}] $-2$。
  \item[\textbf{16.}] $12$。
\end{itemize}

\gkpart{解答题}
\begin{itemize}[leftmargin=2.6em]
  \item[\textbf{17.}] （1）$A=\dfrac{\pi}6$；（2）$\dfrac12$。
  \item[\textbf{18.}] （1）$\dfrac{27}{128}$；（2）$\dfrac{175}{256}$。
  \item[\textbf{19.}] （1）$a=-3$；（2）增区间 $(-\infty,-1)$、$(3,+\infty)$，减区间 $(-1,3)$。
  \item[\textbf{20.}] （1）$\sqrt3$；（2）$\arccos\dfrac{\sqrt6}5$。
  \item[\textbf{21.}] （1）$x^2-\dfrac{y^2}3=1$；（2）$1+\sqrt{17}$。
  \item[\textbf{22.}] （1）$a_3=16$，$a_4=\dfrac1{256}$，$a_{2008}=2^{-2^{2007}}$；（2）$\sqrt2$。
\end{itemize}
